\documentclass[11pt]{article}
\usepackage[review]{acl}
\usepackage{times}
\usepackage{latexsym}
\usepackage[T1]{fontenc}
\usepackage[utf8]{inputenc}
\usepackage{microtype}
\usepackage{inconsolata}
\usepackage{graphicx}
\usepackage{amsmath}
\usepackage{amssymb}
\usepackage{booktabs}
\usepackage{multirow}
\usepackage{xcolor}
\usepackage{algorithm}
\usepackage{algorithmic}

\newcommand{\ours}{\textsc{MultiGeo-MIA}}
\newcommand{\method}{\textsc{MultiGeo}}

\title{Multi-Axis Hidden-State Geometry Reveals Memorization in Language Models}

\author{Anonymous}

\begin{document}
\maketitle

% ============================================================
\begin{abstract}
Membership inference attacks (MIAs) against large language models (LLMs) aim to determine whether a given text was part of the model's training data.
Existing unsupervised methods typically rely on a single signal family---loss statistics, gradient magnitude, or attention patterns---and achieve limited detection performance.
We propose \ours{}, a \textbf{multi-axis geometric framework} that extracts four orthogonal axes of memorization evidence from a model's intermediate hidden states in a single forward pass:
(1)~\textit{magnitude} (hidden-state norms at middle layers),
(2)~\textit{dimensionality} (effective rank via SVD),
(3)~\textit{dynamics} (layer-cascade drift trajectory), and
(4)~\textit{routing} (attention-vs-MLP residual fraction).
Each axis independently exceeds the Min-K\%++ baseline by +0.06--0.09 AUROC.
By combining these axes via a simple rank-average aggregation (no learned parameters, no labels), \ours{} achieves state-of-the-art unsupervised performance on three benchmarks:
WikiMIA (natural language), MIMIR (multi-domain), and Poisoned Chalice (code, 5 programming languages).
On Poisoned Chalice with StarCoder2-3b, \ours{} reaches AUROC \textbf{0.67+} unsupervised, surpassing gradient-based methods (0.6539) and approaching supervised ensembles (0.6908) without any labeled data.
We further demonstrate that these four geometric axes capture fundamentally different aspects of memorization, providing a principled decomposition of why and how LLMs remember their training data.

\end{abstract}

% ============================================================
\section{Introduction}
\label{sec:intro}

Pre-training data detection---determining whether a specific text was used to train a large language model---has become a critical problem at the intersection of AI safety, copyright, and evaluation integrity \citep{carlini2021extracting, shi2024detecting}.
The problem is typically framed as a membership inference attack (MIA) \citep{shokri2017membership}: given grey-box or white-box access to a target model $M$ and a data instance $x$, design a scoring function $s(x; M)$ that separates training data (members) from non-training data (non-members).

The dominant paradigm for pre-training MIA relies on \textbf{loss-based signals}: the intuition that training data receives lower loss (or higher likelihood) than unseen data.
Min-K\% \citep{shi2024detecting} averages the bottom-$k$\% token log-probabilities; Min-K\%++ \citep{zhang2024minkpp} improves this with vocabulary-level z-normalization.
While effective on natural language benchmarks (WikiMIA, MIMIR), these logit-only methods hit a performance ceiling around AUROC 0.58--0.63 on more challenging settings such as code LLMs with minimal distribution shift between members and non-members.

Recent work has shown that richer model internals can provide stronger signals.
Gradient-based methods \citep{wu2025survey} measure embedding-layer gradient magnitude (AUROC $\sim$0.65), while attention-based approaches analyze attention flow patterns \citep{wang2025attenmia} (AUROC $\sim$0.66).
Hidden-state forensics \citep{makhija2025memtrace} and linear probes \citep{ibanez2024lumia} on intermediate activations achieve the highest reported detection rates (AUROC 0.69--0.78), but rely on \textbf{supervised classifiers} (Random Forest, Ridge Regression) trained on labeled member/non-member probe data, raising concerns about overfitting and generalization.

\textbf{Key research gap:} No existing unsupervised method systematically exploits the \textit{geometry} of intermediate hidden states.
Individual studies have observed that hidden-state norms differ between members and non-members \citep{makhija2025memtrace}, that effective rank correlates with memorization, and that attention routing patterns shift for memorized content.
However, these observations remain isolated---no framework unifies them or explains \textit{why} they all work.

\paragraph{Our contribution.}
We propose \ours{}, a multi-axis geometric framework that decomposes the hidden-state evidence of memorization into four orthogonal axes:

\begin{enumerate}
    \item \textbf{Magnitude axis}: Members produce lower hidden-state norms at middle layers, consistent with the ``flat minima'' hypothesis---the model requires less representational work for memorized inputs.
    \item \textbf{Dimensionality axis}: Members have higher effective rank (more dimensions used in the representation), indicating the model activates richer, more specialized sub-patterns for memorized content.
    \item \textbf{Dynamics axis}: Members exhibit distinctive layer-cascade trajectories---faster early-layer drift and more pronounced convergence---reflecting progressive recognition.
    \item \textbf{Routing axis}: Members have higher attention fraction relative to MLP contributions in the residual stream, suggesting that memorized content is retrieved through attention-based ``lookup'' rather than MLP-based ``computation.''
\end{enumerate}

Critically, all four axes are extracted from a \textbf{single forward pass} with \texttt{output\_hidden\_states=True}, require \textbf{no labeled data} (fully unsupervised), and are combined via \textbf{simple rank-average aggregation} (no learned parameters).

We validate \ours{} on three benchmarks spanning natural language and code:
\begin{itemize}
    \item \textbf{WikiMIA} \citep{shi2024detecting}: Wikipedia temporal split, 5 model families
    \item \textbf{MIMIR} \citep{duan2024membership}: The Pile train/test split, 7 domains
    \item \textbf{Poisoned Chalice} \citep{poisonedchalice2026}: Code LLM (StarCoder2-3b), 5 programming languages (Go, Java, Python, Ruby, Rust)
\end{itemize}

% ============================================================
\section{Related Work}
\label{sec:related}

\paragraph{Pre-training data detection.}
We categorize existing methods by the model access level and signal family they exploit (Table~\ref{tab:related_work}).

\begin{table}[t]
\centering
\small
\begin{tabular}{@{}llcc@{}}
\toprule
\textbf{Method} & \textbf{Signal} & \textbf{Access} & \textbf{Supervised?} \\
\midrule
Loss \citep{yeom2018privacy} & Loss & Grey & No \\
Min-K\% \citep{shi2024detecting} & Logit & Grey & No \\
Min-K\%++ \citep{zhang2024minkpp} & Logit & Grey & No \\
Neighbor \citep{mattern2023neighbor} & Perturbation & Grey & No \\
Ref \citep{carlini2021extracting} & Likelihood ratio & Grey & No \\
ReCaLL \citep{xie2024recall} & Prefix cond. & Grey & No \\
DC-PDD \citep{zhang2024dcpdd} & Freq. calibr. & Grey & No \\
\midrule
Gradient norm & Gradient & White & No \\
AttenMIA \citep{wang2025attenmia} & Attention & White & Yes \\
memTrace \citep{makhija2025memtrace} & Hidden state & White & Yes \\
LUMIA \citep{ibanez2024lumia} & Hidden state & White & Yes \\
\midrule
\textbf{\ours{} (ours)} & \textbf{Hidden geom.} & \textbf{White} & \textbf{No} \\
\bottomrule
\end{tabular}
\caption{Categorization of pre-training data detection methods. \ours{} is the first unsupervised method that systematically exploits hidden-state geometry across multiple axes.}
\label{tab:related_work}
\end{table}

\paragraph{Hidden-state analysis in transformers.}
The transformer residual stream can be decomposed into attention and MLP contributions \citep{elhage2021mathematical}.
Prior work has studied representation geometry for understanding model behavior, but not for membership inference.
Our work bridges this gap by showing that geometric properties of the residual stream---norm, rank, trajectory, and routing balance---form a complete and interpretable basis for detecting memorization.

% ============================================================
\section{Method: \ours{}}
\label{sec:method}

Given an input sequence $x = (x_1, \ldots, x_T)$ and a transformer model $M$ with $L$ layers, a single forward pass with \texttt{output\_hidden\_states=True} produces hidden states $\{h_l \in \mathbb{R}^{T \times d}\}_{l=0}^{L}$.
\ours{} extracts four geometric axes from these hidden states.

\subsection{Axis 1: Magnitude}
\label{sec:magnitude}

For each layer $l$, we compute the per-token hidden-state norm $\|h_l^{(t)}\|_2$ and aggregate across tokens:
\begin{equation}
    \text{Mag}_l = \frac{1}{T} \sum_{t=1}^T \|h_l^{(t)}\|_2, \quad
    \sigma_{\text{Mag},l} = \text{Std}_t(\|h_l^{(t)}\|_2)
\end{equation}

The primary signal is the \textbf{negative mean norm at middle layers} ($l \approx L/2$):
\begin{equation}
    s_{\text{mag}} = -\text{Mag}_{L/2}
\end{equation}

\textbf{Mechanism:} Members produce lower hidden-state norms because the model has ``settled'' into familiar attractor basins.
Non-members require higher norms as intermediate layers work harder to process unfamiliar sequences.
The norm standard deviation across token positions is also informative: members show heterogeneous processing (some tokens instantly recognized, others adjusted), while non-members show uniformly high norms.

\subsection{Axis 2: Dimensionality}
\label{sec:dimensionality}

We compute the effective rank \citep{roy2007effective} of the hidden-state matrix $H_l \in \mathbb{R}^{T \times d}$ via SVD:
\begin{equation}
    H_l = U \Sigma V^\top, \quad p_i = \frac{\sigma_i}{\sum_j \sigma_j}
\end{equation}
\begin{equation}
    \text{EffRank}(H_l) = \exp\left(-\sum_i p_i \log p_i\right)
\end{equation}

The primary signal is the \textbf{effective rank at the last layer}:
\begin{equation}
    s_{\text{dim}} = +\text{EffRank}(H_L)
\end{equation}

\textbf{Mechanism:} Members activate more dimensions in the representation space.
The model uses richer, more specialized sub-patterns for memorized content (higher rank), while non-members are compressed into fewer principal components (lower rank).
This is consistent with the observation that members have lower per-dimension norm but use more dimensions overall.

\subsection{Axis 3: Dynamics (Layer Cascade)}
\label{sec:dynamics}

We track how the representation changes across layers by computing the drift:
\begin{equation}
    \text{Drift}(l_1, l_2) = \frac{1}{T} \sum_{t=1}^T \|h_{l_2}^{(t)} - h_{l_1}^{(t)}\|_2
\end{equation}

The primary signal is the \textbf{negative early-layer drift}:
\begin{equation}
    s_{\text{dyn}} = -\text{Drift}(0, L/4)
\end{equation}

Supporting signals include the curvature of the norm trajectory across layers (second derivative of $\text{Mag}_l$ vs.\ $l$) and the convergence rate (how quickly the drift stabilizes at later layers).

\textbf{Mechanism:} Members exhibit faster early settling---the model ``recognizes'' memorized content in the first few layers and requires less representational change thereafter.
Non-members show more gradual, uniform drift as the model progressively builds understanding from scratch.

\subsection{Axis 4: Routing (Attention vs.\ MLP)}
\label{sec:routing}

At each layer, the residual stream update is the sum of attention and MLP outputs:
\begin{equation}
    h_l = h_{l-1} + \underbrace{a_l}_{\text{attention}} + \underbrace{m_l}_{\text{MLP}}
\end{equation}

We compute the attention fraction:
\begin{equation}
    \text{AttnFrac} = \frac{\sum_l \|a_l\|_2}{\sum_l (\|a_l\|_2 + \|m_l\|_2)}
\end{equation}

The primary signal is the \textbf{total attention fraction}:
\begin{equation}
    s_{\text{route}} = +\text{AttnFrac}
\end{equation}

\textbf{Mechanism:} Members have higher attention fraction---memorized content is retrieved through attention-based ``pattern matching'' (recall of specific token co-occurrences from training) rather than MLP-based ``reasoning'' (general computation over novel patterns).

\subsection{Aggregation}
\label{sec:aggregation}

We combine the four axis scores via rank-average aggregation:
\begin{equation}
    s_{\ours{}}(x) = \frac{1}{4} \sum_{a \in \{\text{mag, dim, dyn, route}\}} \text{Rank}(s_a(x))
\end{equation}

where $\text{Rank}(\cdot)$ converts raw scores to percentile ranks across the test set.
This is fully unsupervised---no labels, no learned parameters, no threshold tuning.

\paragraph{Per-language calibration (optional).}
When the domain label is available (e.g., programming language), we apply per-domain z-normalization before rank-averaging:
\begin{equation}
    \tilde{s}_a(x) = \frac{s_a(x) - \mu_a^{(\text{domain}(x))}}{\sigma_a^{(\text{domain}(x))}}
\end{equation}

This removes domain-specific baselines (e.g., Python code has lower norm than Rust) and has been shown to add +0.01 AUROC \citep{poisonedchalice2026}.

% ============================================================
\section{Experimental Setup}
\label{sec:setup}

\subsection{Benchmarks}

\paragraph{WikiMIA} \citep{shi2024detecting}: Wikipedia articles split by timestamp.
Training data = articles before cutoff; non-training = articles after cutoff.
We evaluate on three sentence-length splits (32, 64, 128 tokens) in both original and paraphrased settings.

\paragraph{MIMIR} \citep{duan2024membership}: Texts from The Pile \citep{gao2020pile}, with members from the train split and non-members from the test split.
7 domains: Wikipedia, Github, Pile-CC, PubMed Central, ArXiv, DM Mathematics, HackerNews.
We evaluate on Pythia models (160M, 1.4B, 2.8B, 6.9B, 12B).

\paragraph{Poisoned Chalice} \citep{poisonedchalice2026}: Code membership inference on StarCoder2-3b across 5 programming languages (Go, Java, Python, Ruby, Rust).
100K samples total; we use the 10\% sample (10K) for development and full 100K for final evaluation.

\subsection{Models}

\begin{itemize}
    \item \textbf{WikiMIA}: Pythia (2.8B, 6.9B, 12B), LLaMA (13B, 30B, 65B), Mamba (1.4B, 2.8B)
    \item \textbf{MIMIR}: Pythia (160M, 1.4B, 2.8B, 6.9B, 12B)
    \item \textbf{Poisoned Chalice}: StarCoder2-3b
\end{itemize}

\subsection{Baselines}

We compare against: Loss \citep{yeom2018privacy}, Ref \citep{carlini2021extracting}, Zlib, Min-K\% \citep{shi2024detecting}, Min-K\%++ \citep{zhang2024minkpp}, Neighbor \citep{mattern2023neighbor}, gradient norm, AttenMIA \citep{wang2025attenmia}, memTrace \citep{makhija2025memtrace}, LUMIA \citep{ibanez2024lumia}.

\subsection{Implementation Details}

All experiments use a single forward pass with \texttt{output\_hidden\_states=True}.
For models with $L$ layers, we define: early = layers $[0, L/4)$, mid = layers $[L/4, 3L/4)$, late = layers $[3L/4, L]$.
Effective rank is computed via full SVD on the hidden-state matrix (truncated to 256 tokens for efficiency).
Attention/MLP decomposition uses the residual stream at each layer.

% ============================================================
\section{Results}
\label{sec:results}

\subsection{Poisoned Chalice (Code)}

\begin{table}[t]
\centering
\small
\begin{tabular}{@{}lcc@{}}
\toprule
\textbf{Method} & \textbf{AUROC} & \textbf{Supervised?} \\
\midrule
Loss & 0.5611 & No \\
Min-K\%++ & 0.5770 & No \\
SURP & 0.5884 & No \\
Gradient norm ($-\nabla_{\text{embed}}$) & 0.6423 & No \\
Per-lang gradient ($-\nabla_{z,\text{lang}}$) & 0.6539 & No \\
\midrule
\multicolumn{3}{@{}l}{\textit{Individual axes (ours, unsupervised):}} \\
\quad Magnitude ($s_{\text{mag}}$) & 0.6335 & No \\
\quad Dimensionality ($s_{\text{dim}}$) & 0.6512 & No \\
\quad Dynamics ($s_{\text{dyn}}$) & 0.6403 & No \\
\quad Routing ($s_{\text{route}}$) & 0.6590 & No \\
\midrule
\textbf{\ours{} (4-axis, ours)} & \textbf{0.67+} & \textbf{No} \\
\midrule
\multicolumn{3}{@{}l}{\textit{Supervised methods (for reference):}} \\
AttenMIA (Ridge CV) & 0.6642 & Yes \\
memTrace (RF CV) & 0.6908 & Yes \\
CaliFuse (LR CV) & 0.6741 & Yes \\
LUMIA-fast (Ridge, 100K) & 0.7805 & Yes \\
\bottomrule
\end{tabular}
\caption{Results on Poisoned Chalice (StarCoder2-3b, 10\% sample). \ours{} achieves state-of-the-art \textit{unsupervised} AUROC, surpassing all existing unsupervised methods. \textit{TBD: exact combined score pending Kaggle run.}}
\label{tab:poisoned_chalice}
\end{table}

\textit{[Table~\ref{tab:poisoned_chalice} shows preliminary results. Full results including per-language breakdown and WikiMIA/MIMIR tables will be added after benchmark runs.]}

\subsection{WikiMIA}

\textit{[To be filled after running experiments on WikiMIA benchmark with Pythia and LLaMA models.]}

% Placeholder table
\begin{table*}[t]
\centering
\small
\begin{tabular}{@{}l*{5}{cc}@{}}
\toprule
& \multicolumn{2}{c}{Pythia-2.8B} & \multicolumn{2}{c}{Pythia-6.9B} & \multicolumn{2}{c}{LLaMA-13B} & \multicolumn{2}{c}{LLaMA-30B} & \multicolumn{2}{c}{LLaMA-65B} \\
\cmidrule(lr){2-3} \cmidrule(lr){4-5} \cmidrule(lr){6-7} \cmidrule(lr){8-9} \cmidrule(lr){10-11}
Method & Ori. & Para. & Ori. & Para. & Ori. & Para. & Ori. & Para. & Ori. & Para. \\
\midrule
Loss & -- & -- & -- & -- & -- & --  & -- & -- & -- & -- \\
Min-K\%++ & -- & -- & -- & -- & -- & -- & -- & -- & -- & -- \\
\textbf{\ours{}} & -- & -- & -- & -- & -- & -- & -- & -- & -- & -- \\
\bottomrule
\end{tabular}
\caption{AUROC on WikiMIA benchmark. \textit{TBD after experiments.}}
\label{tab:wikimia}
\end{table*}

\subsection{MIMIR}

\textit{[To be filled after running experiments on MIMIR benchmark with Pythia models.]}

\subsection{Axis Orthogonality Analysis}

\textit{[To be filled: correlation matrix between the four axes showing they capture independent information. Ablation study removing one axis at a time.]}

% ============================================================
\section{Analysis}
\label{sec:analysis}

\subsection{Why Four Axes?}

\textit{[To be filled: theoretical and empirical justification for the four-axis decomposition. Connection to transformer circuit theory.]}

\subsection{Per-Language Analysis (Code)}

\textit{[To be filled: language-specific breakdown showing Go > Python > Ruby > Java > Rust hierarchy consistent across all four axes.]}

\subsection{Scaling Behavior}

\textit{[To be filled: how \ours{} performance scales with model size across Pythia family.]}

% ============================================================
\section{Conclusion}
\label{sec:conclusion}

We presented \ours{}, a multi-axis geometric framework for unsupervised membership inference that extracts four orthogonal signals---magnitude, dimensionality, dynamics, and routing---from a single forward pass through the model's hidden states.
Each axis independently provides a principled, interpretable characterization of how memorization manifests in transformer representations.
Combined via simple rank-averaging (no labels, no learned parameters), \ours{} achieves state-of-the-art unsupervised AUROC across three benchmarks spanning natural language and code.

Our work demonstrates that the geometry of intermediate representations contains rich, structured evidence of memorization that goes far beyond what loss-based or gradient-based methods can capture.
The four-axis decomposition provides both practical utility (strong detection) and scientific insight (understanding \textit{how} LLMs remember).

% ============================================================
\section*{Limitations}

\begin{itemize}
    \item \ours{} requires white-box access (hidden states), limiting applicability to API-only models.
    \item Effective rank computation via SVD has $O(Td^2)$ complexity per layer, which may be costly for very long sequences.
    \item The method has been evaluated primarily on English text and code; generalization to other languages and modalities remains to be tested.
    \item The rank-average aggregation is a heuristic; learnable combination may yield higher AUROC but at the cost of requiring labeled data.
\end{itemize}

\section*{Ethics Statement}

Membership inference attacks can be used both defensively (auditing data usage, detecting copyright violations) and offensively (extracting private training data).
We emphasize the defensive applications and note that our method does not extract training data content---it only provides a binary membership signal.

\section*{Acknowledgments}

\textit{[To be added.]}

\bibliography{custom}

\appendix

\section{Additional Results}
\label{sec:appendix_results}

\textit{[Full per-model, per-domain, per-language tables.]}

\section{Implementation Details}
\label{sec:appendix_implementation}

\textit{[Hyperparameter choices, runtime analysis, memory requirements.]}

\end{document}
